% =============================================================================
% OPERATOR LEARNING PROJECT: HOW TO RUN THE CODE ON A PC AND ON SUNBIRD
% Compile with pdflatex (twice, for the table of contents).
% =============================================================================
\documentclass[11pt,a4paper]{article}

\usepackage[T1]{fontenc}
\usepackage{lmodern}
\usepackage[margin=24mm]{geometry}
\usepackage{xcolor}
\usepackage{listings}
\usepackage{booktabs}
\usepackage{tabularx}
\usepackage{enumitem}
\usepackage{titlesec}
\usepackage{tikz}
\usetikzlibrary{positioning,arrows.meta}
\usepackage[hidelinks]{hyperref}

% COPY-PASTE SAFETY: put real space characters into the PDF, so that text copied
% from a PDF viewer keeps its spaces, and map every glyph to its plain character
\input{glyphtounicode}
\pdfgentounicode=1
\ifdefined\pdfinterwordspaceon \pdfinterwordspaceon \fi

\setlength{\parindent}{0pt}
\setlength{\parskip}{6pt}
\setlist{itemsep=2pt,topsep=2pt}
\titleformat{\part}[display]{\Large\bfseries\centering}{Part \thepart}{4pt}{}
\titleformat{\section}{\large\bfseries}{\thesection}{0.8em}{}
\titleformat{\subsection}{\normalsize\bfseries}{\thesubsection}{0.8em}{}

% INLINE CODE: typed exactly as shown (keeps -- as two hyphens)
\lstset{basicstyle=\ttfamily, columns=fullflexible, keepspaces=true, upquote=true,
        breaklines=true, breakatwhitespace=false, prebreak={}, postbreak={}}
\newcommand{\code}[1]{\lstinline{#1}}
\newcolumntype{Y}{>{\raggedright\arraybackslash}X}

% COMMAND BOXES: copy the text inside exactly as it is
\lstnewenvironment{cmd}[1][]{%
  \lstset{basicstyle=\ttfamily\footnotesize, columns=fullflexible, keepspaces=true,
          upquote=true, showstringspaces=false, breaklines=false,
          frame=single, rulecolor=\color{black!60}, backgroundcolor=\color{black!4},
          xleftmargin=4pt, xrightmargin=4pt, framesep=5pt, aboveskip=6pt, belowskip=6pt,
          title={\footnotesize\sffamily #1}}}{}

% A FRAMED NOTE
\newcommand{\note}[1]{\par\noindent\fbox{\parbox{\dimexpr\linewidth-2\fboxsep-2\fboxrule}{\small #1}}\par}

\begin{document}

% -----------------------------------------------------------------------------
\begin{center}
  {\LARGE\bfseries Operator Learning Project}\\[4pt]
  {\large Running the code on your own computer and on Sunbird}\\[8pt]
  {\small DeepONet and FNO for steady heat conduction \quad$\cdot$\quad October 2026}
\end{center}

\vspace{2pt}
\tableofcontents
\vspace{10pt}

% =============================================================================
\section{The one idea}
% =============================================================================

Everything is done by one program, \code{run.py}. You give it a \emph{stage} (what to do)
and a few \emph{options} (how to do it). The same words work on every machine:

\begin{center}
\begin{tabular}{@{}ll@{}}
\toprule
Where & What you type \\
\midrule
Your computer      & \code{python run.py  STAGE  OPTIONS} \\
Sunbird, GPU node  & \code{sbatch jobs/slurm/gpu.sh  STAGE  OPTIONS} \\
Sunbird, CPU node  & \code{sbatch jobs/slurm/cpu.sh  STAGE  OPTIONS} \\
\bottomrule
\end{tabular}
\end{center}

Example, training an FNO with seed 321:

\begin{cmd}[Your computer]
python run.py train --operator fno --seed 321
\end{cmd}
\begin{cmd}[Sunbird]
sbatch jobs/slurm/gpu.sh train --operator fno --seed 321
\end{cmd}

\note{\textbf{Copying commands.} Every command of this manual is also in the plain-text files
\code{docs/manual/commands_windows.txt} and \code{docs/manual/commands_slurm.txt}.
Open them in Notepad and copy from there: plain text never loses a space.}

\note{\textbf{Spaces matter.} Every word is separated by one space, and every option starts
with two hyphens: \code{smoke --device cpu}, never \code{smoke--device cpu}.
If the computer answers \code{invalid choice}, look for a missing space.}

See Section~\ref{sec:order} for the whole study in the order you run it.

\textbf{Where things are.} Everything is inside the project folder (the folder that contains
\code{run.py}). The code creates two folders: \code{data/} for the dataset and
\code{outputs/} for every result. A \emph{run} is one trained model with one name, for example
\code{final_fno_seed321}. Its results are:

\begin{center}
\begin{tabularx}{\textwidth}{@{}lY@{}}
\toprule
Folder & Contents \\
\midrule
\code{outputs/checkpoints/combined_0500/RUN/}   & trained weights \code{best_model.pt} \\
\code{outputs/logs/combined_0500/RUN/}          & loss history, training summary \\
\code{outputs/study_results/combined_0500/RUN/} & error tables (validation, test) \\
\code{outputs/figures/combined_0500/RUN/}       & figures of this run \\
\midrule
\code{outputs/figures/combined_0500/report/}    & tuning heatmaps and DeepONet--FNO comparison \\
\code{outputs/figures/combined_0500/dataset/}   & where the cavities lie \\
\bottomrule
\end{tabularx}
\end{center}

\code{combined_0500} is the earlier dataset of 500 cases (circles and rectangles). The runs of the
full study (Section~\ref{sec:order}) are in the same folders under \code{scaling_3200} instead.

% =============================================================================
\section{The full study, in order}\label{sec:order}
% =============================================================================

\textbf{The question.} How accurate are DeepONet and FNO, and how does each improve when it is
given more training data? The study therefore has a hyperparameter of its own: the number of
training cases $N \in \{200, 400, 800, 1600, 3200\}$.

\textbf{What ``scaling'' means.} It is the usual name for \emph{how the test error scales with
$N$}. For a well-behaved model the error falls like a power of $N$,
\[
  e(N) = C\,N^{-\alpha}
  \qquad\Longleftrightarrow\qquad
  \log e = \log C - \alpha \log N ,
\]
a straight line on log--log axes. The slope $\alpha$ is the \emph{rate}: doubling the data
multiplies the error by $2^{-\alpha}$ (for $\alpha = 0.5$, by 0.71; for $\alpha = 1$, by 0.5).
A model with the larger $\alpha$ profits more from data; a curve that bends flat means the model
has stopped learning from data (its capacity, not the data, is the limit).

\textbf{One dataset for every $N$} (\code{scaling_3200}, 3800 FE solves):
\begin{center}
\begin{tabularx}{\textwidth}{@{}llY@{}}
\toprule
Part & Cases & Role \\
\midrule
training pool & 3200 & Halton sequence, alternating circle/rectangle. The training set of size $N$
                        is the first $N$ cases: each set lies inside the next and is evenly spread. \\
validation    & 200  & Latin hypercube. \textbf{The same 200 cases for every $N$ and every seed.}
                        Used for early stopping, the learning-rate schedule and tuning. \\
test          & 400  & Latin hypercube. \textbf{The same 400 cases for every $N$ and every seed.}
                        Used only for the final errors. \\
\bottomrule
\end{tabularx}
\end{center}
Keeping validation and test fixed is essential: every point of the curve is then measured on
identical cases, so a change of error is caused by $N$ alone. The temperature scale is taken from
the whole training pool, so it is also the same for every $N$.

\textbf{Every run is trained the same way:} at most 5000 epochs, early stopping after 800 epochs
without a better validation error, learning rate halved after 200 such epochs, batch size 8, the
tuned architecture. Each $N$ is trained with five seeds (321, 777, 2026, 31415, 27182): 25 runs per
operator, 50 in total. The seeds give the spread (error bars) of each point.

\textbf{The order.} Each step needs the results of the one before.

\begin{center}
\begin{tabularx}{\textwidth}{@{}clYl@{}}
\toprule
Step & Stage & What happens & Section \\
\midrule
1 & \code{generate}, \code{generate-rounded} & FE solves: the 3800 cases, and 50 rounded rectangles
                                               (a shape never trained on) & \ref{sec:pcdata} \\
2 & \code{plot-sampling}, \code{preview}, \code{mesh-check} & check the data: where the cavities lie,
                                               one case drawn, FE accuracy & \ref{sec:pcdata} \\
3 & \code{train} (short) & optional: one quick model, to see that training works & \ref{sec:pcone} \\
4 & \code{tune}, \code{tune-select} & choose each operator's architecture and learning rate at
                                               $N = 800$, on the validation cases only & \ref{sec:pctune} \\
5 & \code{train-scaling} & the 25 runs per operator, each scored on the test cases & \ref{sec:pcscaling} \\
6 & \code{scaling-report} & error against $N$, fitted $\alpha$ & \ref{sec:pcscaling} \\
7 & \code{compare} & full DeepONet--FNO report at each $N$ & \ref{sec:pcscaling} \\
8 & \code{plot}, \code{evaluate-rounded}, \code{benchmark} & figures, rounded-rectangle errors and
                                               timings of chosen runs & \ref{sec:pcscaling} \\
\bottomrule
\end{tabularx}
\end{center}

Why tune at $N = 800$: tuning at every $N$ would cost five times as much. 800 is in the middle of
the range (on log axes), so the chosen architecture is not tailored to the smallest or the largest
set. Every $N$ then uses the same architecture, so only $N$ changes along the curve.

\textbf{The earlier 500-case study} (\code{combined_0500}: one Latin-hypercube dataset, split
400/50/50) uses the same stages without \code{--family scaling}: tuning, then \code{train-final}
with five seeds, then \code{compare}. It is kept to reproduce the earlier results and is not needed
for the full study.

\textbf{What you can change.} Everything has a default; you type an option only to change it.
Section~\ref{sec:options} lists every option of \code{run.py}, and
Section~\ref{sec:settings} the settings that are fixed in \code{src/config.py} (physics, geometry
ranges, grid). Options needed for the full study are few: \code{--family scaling},
\code{--operator}, \code{--stage-name}, \code{--train-size 800} for tuning, and
\code{--epochs 2000} for tuning.

% =============================================================================
\part{Your own computer (Windows, NVIDIA GPU)}
% =============================================================================

Do Sections~\ref{sec:driver}--\ref{sec:project} once. After that, every session starts with
the two lines of Section~\ref{sec:everyday}.

% =============================================================================
\section{NVIDIA driver}\label{sec:driver}
% =============================================================================

PyTorch brings its own copy of CUDA. You do \textbf{not} install the CUDA Toolkit;
you only need a recent NVIDIA \emph{driver}.

\begin{enumerate}
  \item Press the Windows key, type \code{cmd}, press Enter. Type:
\begin{cmd}[Command Prompt]
nvidia-smi
\end{cmd}
  \item A table appears with the name of your graphics card. In its top-right corner it shows
        \code{CUDA Version: 13.0} (or another number). Write this number down.
  \item If the command is not found, or the number is below 12.6, install the newest driver:
        go to \href{https://www.nvidia.com/Download/index.aspx}{nvidia.com/Download}, choose your
        card, download, run the installer with the default options, restart the computer, and
        repeat step~1.
\end{enumerate}

% =============================================================================
\section{Miniconda}\label{sec:conda}
% =============================================================================

Miniconda gives you Python in separate, clean \emph{environments}.

\begin{enumerate}
  \item Download \emph{Miniconda3 Windows 64-bit} from
        \href{https://www.anaconda.com/download/success}{anaconda.com/download/success}
        (section \emph{Miniconda Installers}).
  \item Run the installer with these choices:
        \emph{Just Me}; keep the suggested folder
        (\texttt{miniconda3} in your user folder); leave \emph{Add to PATH} unticked. Finish.
  \item Open the Start menu and start \textbf{Anaconda Prompt (miniconda3)}.
        Every command in Part~I is typed in this window. Its prompt starts with \code{(base)}.
\end{enumerate}

% =============================================================================
\section{Python environment}\label{sec:env}
% =============================================================================

Create an environment called \code{opl}. The options \code{--override-channels -c conda-forge}
take Python from the free conda-forge channel and avoid the error
\emph{Terms of Service have not been accepted}.

\begin{cmd}[Anaconda Prompt]
conda create -n opl python=3.11 -y --override-channels -c conda-forge
conda activate opl
pip install numpy matplotlib ngsolve
\end{cmd}

The prompt now starts with \code{(opl)}.

\textbf{PyTorch with GPU.} Use the box that matches the CUDA Version you wrote down in
Section~\ref{sec:driver}:

\begin{cmd}[nvidia-smi shows CUDA Version 13.0 or higher]
pip install torch --index-url https://download.pytorch.org/whl/cu130
\end{cmd}
\begin{cmd}[nvidia-smi shows CUDA Version 12.6 to 12.9]
pip install torch --index-url https://download.pytorch.org/whl/cu126
\end{cmd}

If PyTorch was installed earlier with a plain \code{pip install torch}, it is the CPU version.
Remove it first:

\begin{cmd}[Anaconda Prompt]
pip uninstall -y torch
pip install torch --index-url https://download.pytorch.org/whl/cu130
\end{cmd}

\textbf{Check the GPU.}

\begin{cmd}[Anaconda Prompt]
python -c "import torch; print(torch.__version__, torch.cuda.is_available())"
python -c "import torch; print(torch.cuda.get_device_name(0))"
\end{cmd}

The first line must end with \code{True} and the version with \code{+cu130} (or \code{+cu126});
the second prints the name of your card. If it prints \code{False}, the CPU version is
installed or the driver is too old: repeat this section.
The exact command for other cases is on
\href{https://pytorch.org/get-started/locally/}{pytorch.org/get-started/locally}
(choose Stable, Windows, Pip, Python, CUDA).

% =============================================================================
\section{The project}\label{sec:project}
% =============================================================================

Download the project into \texttt{C:\textbackslash ML}, either with git (Git for Windows,
\code{git-scm.com}):

\begin{cmd}[Anaconda Prompt]
git clone https://github.com/prajwal1798/Feature_Extraction.git C:\ML
\end{cmd}

or from the GitHub page (\emph{Code}, \emph{Download ZIP}): extract it and rename the folder to
\texttt{C:\textbackslash ML}, so that \code{run.py} lies directly inside it. Then:

\begin{cmd}[Anaconda Prompt]
cd /d C:\ML
python tests\run_all.py
python run.py smoke
\end{cmd}

\code{run_all.py} checks the numerics (about one minute) and ends with \emph{All tests passed}.
\code{smoke} runs every stage on 20 tiny cases on your GPU (a few minutes) and ends with
\emph{Smoke test passed}. It checks that everything runs, not the accuracy. Its results go to
the folder \code{verification}, which you may delete afterwards.

% =============================================================================
\section{Every new session}\label{sec:everyday}
% =============================================================================

Open \textbf{Anaconda Prompt (miniconda3)} and type:

\begin{cmd}[Anaconda Prompt]
conda activate opl
cd /d C:\ML
\end{cmd}

Then type the commands of the following sections. To look at results, open the folder in
Explorer with \code{explorer outputs}. Figures are PDF (for the thesis) and PNG (for slides).

\note{\textbf{Long runs.} A training run takes hours. Keep the window open, and set Windows
\emph{Settings > System > Power > Sleep} to \emph{Never} while it runs.
\code{Ctrl+C} stops a run; Section~\ref{sec:resume} shows how to continue it.}

% =============================================================================
\section{Step 1--2: data}\label{sec:pcdata}
% =============================================================================

The finite element solver runs on the processor. 3800 cases take hours; if the command stops
(or you close the window), type it again: it continues where it stopped.

\begin{cmd}[Anaconda Prompt]
python run.py generate --family scaling
python run.py generate-rounded --family scaling --rounded-cases 50
\end{cmd}

Then check the data:

\begin{cmd}[Anaconda Prompt]
python run.py plot-sampling --family scaling
python run.py preview
python run.py mesh-check
\end{cmd}

\begin{center}
\begin{tabularx}{\textwidth}{@{}lY@{}}
\toprule
Stage & What you get \\
\midrule
\code{generate}         & \code{data/processed/scaling_3200/protocol_v2/dataset.npz} (all cases, split, scale) \\
\code{generate-rounded} & 50 rounded rectangles, a shape the networks never see in training \\
\code{plot-sampling}    & where the cavities lie: \code{outputs/figures/scaling_3200/dataset/} \\
\code{preview}          & one FE case (geometry, mesh, temperature, SDF) \\
\code{mesh-check}       & FE error on three meshes and the heat balance (table) \\
\bottomrule
\end{tabularx}
\end{center}

% =============================================================================
\section{Step 3: train, test and plot one model}\label{sec:pcone}
% =============================================================================

Optional, but a good first try. \code{--train-size 800} uses the first 800 training cases;
\code{--epochs 300} keeps it short.

\begin{cmd}[Anaconda Prompt]
set F=--family scaling --operator fno
python run.py train %F% --train-size 800 --seed 321 --run-name try_fno --epochs 300
\end{cmd}

(\code{set F=...} stores the common words once; \texttt{\%F\%} pastes them back.)
The window prints one line every 20 epochs:

\begin{cmd}[What a training line means]
Epoch   400 | training MSE 1.2e-05 | validation MSE 1.9e-05 | RMSE 0.98 K
\end{cmd}

RMSE is the validation error in kelvin. Training stops at the epoch limit, or earlier when the
validation error has not improved for 800 epochs. The weights with the lowest validation error
are kept. When it has finished:

\begin{cmd}[Anaconda Prompt]
python run.py evaluate         %F% --run-name try_fno
python run.py evaluate-rounded %F% --run-name try_fno
python run.py plot             %F% --run-name try_fno
python run.py benchmark        %F% --run-name try_fno
\end{cmd}

\begin{center}
\begin{tabularx}{\textwidth}{@{}lY@{}}
\toprule
Stage & What you get \\
\midrule
\code{evaluate}         & test errors: \code{outputs/study_results/scaling_3200/try_fno/all_test_cases_grid_*} \\
\code{evaluate-rounded} & errors on the rounded rectangles (same folder, subfolder \code{rounded_rect_external_test}) \\
\code{plot}             & loss curve and the first test case (FE, prediction, error):
                          \code{outputs/figures/scaling_3200/try_fno/} \\
\code{benchmark}        & training and prediction times: \code{.../try_fno/benchmarks/} \\
\bottomrule
\end{tabularx}
\end{center}

\textbf{Another test case:} add \code{--test-case 7} (the 400 test cases are numbered 0 to 399).
\textbf{DeepONet:} \code{set F=--family scaling --operator deeponet} and another run name.
\textbf{Other architectures:} add \code{--variant NAME} to \code{train}:

\begin{center}
\begin{tabularx}{\textwidth}{@{}lY@{}}
\toprule
Operator & Variants (default \code{baseline}) \\
\midrule
\code{deeponet} & \code{baseline} (2 hidden layers, tanh), \code{activation_gelu}, \code{depth3},
                  \code{latent256}, \code{depth4_gelu}, \code{depth6_gelu} \\
\code{fno}      & \code{baseline} (4 blocks, width 32, modes $12\times16$),
                  \code{activation_tanh}, \code{depth3}, \code{modes8} \\
\bottomrule
\end{tabularx}
\end{center}

Give every experiment its own \code{--run-name}; the code refuses to overwrite an existing run.

% =============================================================================
\section{Step 4: tuning at N = 800}\label{sec:pctune}
% =============================================================================

Tuning trains many models on the first 800 training cases and keeps the one with the smallest
\emph{validation} error; the test cases are never used. It goes in stages, each starting from the
winner of the previous one, so keep the order:

\begin{center}
\begin{tabularx}{\textwidth}{@{}llYr@{}}
\toprule
Operator & Stage & What changes & Models \\
\midrule
DeepONet & \code{grid}        & layers $\{2,4,6,8,10\}$ $\times$ neurons $\{10,20,40,80,160\}$, tanh & 25 \\
DeepONet & \code{grid_gelu}   & the same 25, GELU                                              & 25 \\
DeepONet & \code{lr}          & learning rate $\{3\cdot10^{-4},\,10^{-3},\,3\cdot10^{-3}\}$ at the better grid winner & 3 \\
FNO      & \code{modes_width} & modes $\{4,8,12,16\}$ $\times$ width $\{16,32,64\}$              & 12 \\
FNO      & \code{depth}       & Fourier blocks $\{2,3,4,6\}$                                   & 4 \\
FNO      & \code{lr}          & learning rate, as above                                        & 3 \\
\bottomrule
\end{tabularx}
\end{center}

Each stage is two commands: \code{tune} trains all its models one after another (2000 epochs
each, the same for all), \code{tune-select} picks the winner and draws the heatmap.

\begin{cmd}[Anaconda Prompt: DeepONet]
set T=--family scaling --train-size 800 --operator deeponet
python run.py tune        %T% --stage-name grid --epochs 2000
python run.py tune-select %T% --stage-name grid
python run.py tune        %T% --stage-name grid_gelu --epochs 2000
python run.py tune-select %T% --stage-name grid_gelu
python run.py tune        %T% --stage-name lr --epochs 2000
python run.py tune-select %T% --stage-name lr
\end{cmd}

\begin{cmd}[Anaconda Prompt: FNO]
set T=--family scaling --train-size 800 --operator fno
python run.py tune        %T% --stage-name modes_width --epochs 2000
python run.py tune-select %T% --stage-name modes_width
python run.py tune        %T% --stage-name depth --epochs 2000
python run.py tune-select %T% --stage-name depth
python run.py tune        %T% --stage-name lr --epochs 2000
python run.py tune-select %T% --stage-name lr
\end{cmd}

If a \code{tune} command is stopped, type it again: finished models are skipped and the
interrupted one continues. Heatmaps (error in every cell, winner framed in red):
\code{outputs/figures/scaling_3200/report/N0800/tuning_*.pdf}; numbers:
\code{outputs/study_results/scaling_3200/tuning/*.csv}.

% =============================================================================
\section{Steps 5--8: the data-size runs and the results}\label{sec:pcscaling}
% =============================================================================

\textbf{Step 5.} One command per operator trains every size $N$ with every seed (25 runs), one
after another, with the winner of the last tuning stage, and scores each run on the 400 test
cases. Runs are named \code{scale_fno_N0800_seed321} and so on. If stopped, type the same command
again: finished runs are skipped and an interrupted one continues.

\begin{cmd}[Anaconda Prompt]
python run.py train-scaling --operator deeponet
python run.py train-scaling --operator fno
\end{cmd}

\textbf{Steps 6 and 7.} The learning curves, then the full DeepONet--FNO report at every $N$
(the \code{for} line repeats the command for every size \texttt{\%n}):

\begin{cmd}[Anaconda Prompt]
python run.py scaling-report
for %n in (200 400 800 1600 3200) do python run.py compare --family scaling --train-size %n
\end{cmd}

\begin{center}
\begin{tabularx}{\textwidth}{@{}lY@{}}
\toprule
Stage & What you get (in \code{outputs/figures/scaling_3200/report/}) \\
\midrule
\code{scaling-report} & \code{scaling.pdf}: test error against $N$ on log--log axes, mean and
                        spread over seeds, fitted line and $\alpha$ for each operator;
                        numbers in \code{scaling.csv} and \code{scaling_rates.json} \\
\code{compare}        & \code{N0200/} \ldots\ \code{N3200/}: learning curves, test error over seeds,
                        case errors, error per shape, DeepONet against FNO case by case \\
\bottomrule
\end{tabularx}
\end{center}

\textbf{Step 8.} Figures, rounded-rectangle errors and timings of single runs, for example
$N = 3200$, seed 321:

\begin{cmd}[Anaconda Prompt]
set R=--family scaling --operator fno --run-name scale_fno_N3200_seed321
python run.py plot             %R%
python run.py evaluate-rounded %R%
python run.py benchmark        %R%
for /L %i in (0,1,399) do python run.py plot %R% --test-case %i
\end{cmd}

The last line draws every one of the 400 test cases. For DeepONet, change both \code{fno} to
\code{deeponet}.

\textbf{Smaller study.} \code{--sizes 200,400,800} and \code{--seeds 321,777} can be added to
\code{train-scaling}, \code{scaling-report} and \code{compare} (the same lists to all three).
One single run: \code{train-scaling --operator fno --train-size 400 --seed 321}.

% =============================================================================
\section{The earlier 500-case study}\label{sec:pcfinal}
% =============================================================================

This reproduces the earlier results on \code{combined_0500}. Its commands carry no
\code{--family}: \code{combined} with 500 cases is the default. Compute the dataset (the same
code and seed as before give the same 500 cases):

\begin{cmd}[Anaconda Prompt]
python run.py generate --cases 500
python run.py generate-rounded --rounded-cases 50
\end{cmd}

or, if it was made on Sunbird, copy it from there (off campus, connect to the university VPN
first; each \code{scp} asks for your Sunbird password, which is not shown while you type):

\begin{cmd}[Anaconda Prompt]
set SUNBIRD=USERNAME@LOGIN_NODE:/scratch/USERNAME/ML
set DATA=data/processed/combined_0500/protocol_v2
mkdir data\processed\combined_0500\protocol_v2
scp %SUNBIRD%/%DATA%/dataset.npz %DATA%/
scp %SUNBIRD%/%DATA%/dataset_manifest.json %DATA%/
python run.py generate-rounded --rounded-cases 50
\end{cmd}

Tune with the commands of Section~\ref{sec:pctune}, but with
\code{set T=--operator deeponet} and \code{set T=--operator fno}. Then train the winner with
five seeds, test, time, and report:

\begin{cmd}[Anaconda Prompt]
set S=321 777 2026 31415 27182
for %s in (%S%) do python run.py train-final --operator deeponet --seed %s
for %s in (%S%) do python run.py train-final --operator fno --seed %s
for %s in (%S%) do python run.py evaluate --operator deeponet --run-name final_deeponet_seed%s
for %s in (%S%) do python run.py evaluate --operator fno --run-name final_fno_seed%s
for %s in (%S%) do python run.py benchmark --operator deeponet --run-name final_deeponet_seed%s
for %s in (%S%) do python run.py benchmark --operator fno --run-name final_fno_seed%s
python run.py compare
\end{cmd}

The report is in \code{outputs/figures/combined_0500/report/}.

% =============================================================================
\section{Stopping and continuing}\label{sec:resume}
% =============================================================================

The code saves a recovery point every 50 epochs.
\begin{itemize}
  \item \code{train} and \code{train-final}: type the same command again with \code{--resume}
        at the end.
  \item \code{tune}, \code{train-scaling} and \code{generate}: type the same command again.
\end{itemize}
Do not change other settings when continuing; the code checks them.

% =============================================================================
\part{Sunbird}
% =============================================================================

The commands are those of Part~I, with \code{python run.py} replaced by
\code{sbatch jobs/slurm/gpu.sh} (networks) or \code{sbatch jobs/slurm/cpu.sh}
(everything else). Jobs run without you; you only submit and watch.

% =============================================================================
\section{Set up once}
% =============================================================================

This part is written for Sunbird (Supercomputing Wales) and works on any Slurm cluster.
\code{USERNAME} and \code{LOGIN_NODE} stand for your user name and the login address of your
cluster; if its work disk is not \code{/scratch/\$USER}, use its own path throughout.

\textbf{Log in} from your computer:

\begin{cmd}[Your computer]
ssh USERNAME@LOGIN_NODE
\end{cmd}

Everything of the project lives in one folder, \code{/scratch/\$USER/ML}: the code, the Python
environment (\code{ML/env}), the data (\code{ML/data}) and every result (\code{ML/outputs}).
Download the project into it:

\begin{cmd}[Cluster]
cd /scratch/$USER
git clone https://github.com/prajwal1798/Feature_Extraction.git ML
cd ML
mkdir -p outputs/slurm
\end{cmd}

\textbf{The job files.} Open \code{jobs/slurm/gpu.sh} and \code{jobs/slurm/cpu.sh} (for example
with \code{nano}) and set the lines marked at their top: \code{--account} (your project code),
\code{--partition} and, for the GPU job, \code{--gres} (one GPU, type name as on your cluster).

\textbf{Python environment} (once; the job files use \code{ML/env/bin/python}; the name of
the conda module differs between clusters, see \code{module avail conda}):

\begin{cmd}[Sunbird]
module load anaconda/2024.06
source "$(conda info --base)/etc/profile.d/conda.sh"
conda create --prefix /scratch/$USER/ML/env python=3.11 -y \
             --override-channels -c conda-forge
conda activate /scratch/$USER/ML/env
/scratch/$USER/ML/env/bin/python -m pip install --no-cache-dir --only-binary=:all: \
       numpy matplotlib ngsolve torch
\end{cmd}

\textbf{Every later login:} \code{cd /scratch/\$USER/ML}. Every \code{sbatch} line must be
typed in this folder.

\textbf{Check.}

\begin{cmd}[Sunbird]
sbatch jobs/slurm/gpu.sh gpu-info
sbatch jobs/slurm/cpu.sh smoke
\end{cmd}

Watch them live (NUMBER is what \code{sbatch} printed; Section~\ref{sec:watch}):

\begin{cmd}[Sunbird]
tail -F outputs/slurm/gpu_NUMBER.out
tail -F outputs/slurm/cpu_NUMBER.out
\end{cmd}
 The GPU report must name an A100; the smoke test must end with
\emph{Smoke test passed} and \code{EXIT CODE: 0}.
The account, partitions and the 24-hour limit are in the top lines of
\code{gpu.sh} and \code{cpu.sh}.

% =============================================================================
\section{Watching jobs and finding results}\label{sec:watch}
% =============================================================================

\textbf{The output file.} \code{sbatch} answers \code{Submitted batch job 8206449}; 8206449 is
the job number. Everything the job prints goes into one file:

\begin{center}
\begin{tabular}{@{}ll@{}}
\toprule
Sent with & Output file \\
\midrule
\code{gpu.sh} & \code{outputs/slurm/gpu_8206449.out} \\
\code{cpu.sh} & \code{outputs/slurm/cpu_8206449.out} \\
\bottomrule
\end{tabular}
\end{center}

It starts with \code{JOB}, \code{NODE}, \code{STARTED} and \code{COMMAND}, and ends with
\code{FINISHED} and \code{EXIT CODE} (0 means success; otherwise the error is just above).

\begin{center}
\small
\begin{tabularx}{\textwidth}{@{}lY@{}}
\toprule
Command & Meaning \\
\midrule
\code{squeue -u $USER}                        & my jobs: \code{PD} waiting, \code{R} running,
                                                \code{CG} finishing; not listed = over \\
\code{tail -F outputs/slurm/cpu_8206449.out}  & the output \textbf{live}, line by line; waits if the
                                                job has not started. \code{Ctrl+C} stops
                                                watching, not the job \\
\code{cat outputs/slurm/cpu_8206449.out}      & the whole output of a finished job \\
\code{tail -n 30 outputs/slurm/cpu_8206449.out} & its last 30 lines (errors are at the end) \\
\code{ls -t outputs/slurm | head}             & newest output files first \\
\code{sacct -j 8206449}                       & \code{COMPLETED}, \code{FAILED}, \code{TIMEOUT} or
                                                \code{CANCELLED} \\
\code{scancel 8206449}                        & stop the job \\
\bottomrule
\end{tabularx}
\end{center}

\textbf{Array jobs.} Every array line of this manual adds
\code{--output=outputs/slurm/gpu_\%A_\%a.out} (stored once as \code{A="..."} and pasted as
\code{\$A}), so task 3 of array 8206500 (shown by
\code{squeue} as \code{8206500_3}) writes \code{outputs/slurm/gpu_8206500_3.out}. Its GPU:
\code{srun --jobid=8206500_3 --overlap nvidia-smi}. A user may hold a limited number of GPUs at
once; further tasks wait with the reason \code{QOSMaxGRESPerUser} and start automatically.

\textbf{The results} are not in \code{outputs/slurm}. Where each stage writes them
(\code{TAG} is \code{scaling_3200} for the full study, \code{combined_0500} for the 500-case
study; \code{RUN} is the run name):

\begin{center}
\small
\begin{tabularx}{\textwidth}{@{}lY@{}}
\toprule
Stage & Results \\
\midrule
\code{gpu-info}         & in the output file; also \code{outputs/study_results/combined_0500/hardware_preflight.json} \\
\code{smoke}            & output file ends with \emph{Smoke test passed}; files in \code{verification/} \\
\code{generate}         & \code{data/processed/TAG/protocol_v2/dataset.npz} \\
\code{generate-rounded} & \code{data/external/TAG/protocol_v2/rounded_rect_external_test/} \\
\code{plot-sampling}    & \code{outputs/figures/TAG/dataset/} \\
\code{mesh-check}       & \code{outputs/study_results/combined_0500/mesh_convergence/} \\
\code{train}, \code{train-final}, \code{train-scaling}
                        & model \code{outputs/checkpoints/TAG/RUN/best_model.pt};
                          loss history \code{outputs/logs/TAG/RUN/}; test errors
                          \code{outputs/study_results/TAG/RUN/} \\
\code{evaluate-rounded} & \code{outputs/study_results/TAG/RUN/rounded_rect_external_test/} \\
\code{plot}             & \code{outputs/figures/TAG/RUN/} \\
\code{benchmark}        & \code{outputs/study_results/TAG/RUN/benchmarks/} \\
\code{tune-select}      & tables and winner \code{outputs/study_results/TAG/tuning/};
                          heatmaps \code{outputs/figures/TAG/report/} (scaling: \code{report/N0800/}) \\
\code{scaling-report}   & \code{outputs/figures/scaling_3200/report/scaling.pdf} (+ \code{.csv}, \code{.json}) \\
\code{compare}          & \code{outputs/figures/TAG/report/} (scaling: subfolder \code{N0800} etc.) \\
\bottomrule
\end{tabularx}
\end{center}

Jobs keep running after you log out. To look at figures, bring \code{outputs} home
(end of Section~\ref{sec:sbstudy}).

% =============================================================================
\section{The study on Sunbird}\label{sec:sbstudy}
% =============================================================================

The same steps as Section~\ref{sec:order}. Submit one block, wait until \code{squeue -u \$USER}
is empty, then submit the next. A line ending in \texttt{\textbackslash} continues on the next
line; copy both together. \code{--array=0-24} starts tasks 0 to 24 at the same time; each task
trains one model. After every \code{sbatch}, follow it live with
\code{tail -F outputs/slurm/gpu_NUMBER.out} (or \code{cpu_}); where the results go is in
Section~\ref{sec:watch}.

\textbf{Steps 1--2: data} (CPU node; 3800 FE solves; if the 24-hour limit ends it, submit again):

\begin{cmd}[Sunbird]
sbatch jobs/slurm/cpu.sh generate --family scaling
\end{cmd}
\begin{cmd}[Sunbird: when it has finished]
sbatch jobs/slurm/cpu.sh generate-rounded --family scaling --rounded-cases 50
sbatch jobs/slurm/cpu.sh plot-sampling --family scaling
sbatch jobs/slurm/cpu.sh mesh-check
\end{cmd}

\textbf{Step 3: one quick model} (optional):

\begin{cmd}[Sunbird]
F="--family scaling --operator fno"
sbatch jobs/slurm/gpu.sh train $F --train-size 800 --seed 321 --run-name try_fno --epochs 300
\end{cmd}

\textbf{Step 4: tuning at $N = 800$.} One array per stage; select when the array has finished.

\begin{cmd}[Sunbird: DeepONet]
T="--family scaling --train-size 800 --operator deeponet"
A="--output=outputs/slurm/gpu_%A_%a.out"
sbatch --array=0-24 $A jobs/slurm/gpu.sh tune $T --stage-name grid --epochs 2000
sbatch jobs/slurm/cpu.sh tune-select $T --stage-name grid
sbatch --array=0-24 $A jobs/slurm/gpu.sh tune $T --stage-name grid_gelu --epochs 2000
sbatch jobs/slurm/cpu.sh tune-select $T --stage-name grid_gelu
sbatch --array=0-2 $A jobs/slurm/gpu.sh tune $T --stage-name lr --epochs 2000
sbatch jobs/slurm/cpu.sh tune-select $T --stage-name lr
\end{cmd}

\begin{cmd}[Sunbird: FNO]
T="--family scaling --train-size 800 --operator fno"
A="--output=outputs/slurm/gpu_%A_%a.out"
sbatch --array=0-11 $A jobs/slurm/gpu.sh tune $T --stage-name modes_width --epochs 2000
sbatch jobs/slurm/cpu.sh tune-select $T --stage-name modes_width
sbatch --array=0-3 $A jobs/slurm/gpu.sh tune $T --stage-name depth --epochs 2000
sbatch jobs/slurm/cpu.sh tune-select $T --stage-name depth
sbatch --array=0-2 $A jobs/slurm/gpu.sh tune $T --stage-name lr --epochs 2000
sbatch jobs/slurm/cpu.sh tune-select $T --stage-name lr
\end{cmd}

(\code{T="..."} stores the common words; \code{\$T} pastes them back. Set \code{T} again in a new
login.)

\textbf{Step 5: the data-size runs.} 5 sizes $\times$ 5 seeds = 25 tasks per operator
(task $k$ is size $k \div 5$, seed $k \bmod 5$):

\begin{cmd}[Sunbird]
A="--output=outputs/slurm/gpu_%A_%a.out"
sbatch --array=0-24 $A jobs/slurm/gpu.sh train-scaling --operator deeponet
sbatch --array=0-24 $A jobs/slurm/gpu.sh train-scaling --operator fno
\end{cmd}

Large training sets may need more than the 24-hour limit. Submit the same lines again: finished
runs are skipped and interrupted ones continue from their recovery point.

\textbf{Steps 6--8: results.}

\begin{cmd}[Sunbird: when all runs have finished]
sbatch jobs/slurm/cpu.sh scaling-report
for n in 200 400 800 1600 3200
do
    sbatch jobs/slurm/cpu.sh compare --family scaling --train-size $n
done
R="--family scaling --operator fno --run-name scale_fno_N3200_seed321"
sbatch jobs/slurm/gpu.sh plot $R
sbatch jobs/slurm/gpu.sh evaluate-rounded $R
sbatch jobs/slurm/gpu.sh benchmark $R
\end{cmd}

\textbf{The earlier 500-case study} (Section~\ref{sec:pcfinal}). Generate its dataset, then
tune with \code{T="--operator deeponet"} and \code{T="--operator fno"} in the arrays above, then
train the final seeds:

\begin{cmd}[Sunbird]
sbatch jobs/slurm/cpu.sh generate --cases 500
\end{cmd}
\begin{cmd}[Sunbird: when it has finished]
sbatch jobs/slurm/cpu.sh generate-rounded --rounded-cases 50
\end{cmd}

\begin{cmd}[Sunbird: after tuning]
for s in 321 777 2026 31415 27182
do
    sbatch jobs/slurm/gpu.sh train-final --operator deeponet --seed $s
    sbatch jobs/slurm/gpu.sh train-final --operator fno      --seed $s
done
\end{cmd}

\begin{cmd}[Sunbird: when the training jobs have finished]
for s in 321 777 2026 31415 27182
do
    sbatch jobs/slurm/gpu.sh evaluate  --operator deeponet --run-name final_deeponet_seed$s
    sbatch jobs/slurm/gpu.sh evaluate  --operator fno      --run-name final_fno_seed$s
    sbatch jobs/slurm/gpu.sh benchmark --operator deeponet --run-name final_deeponet_seed$s
    sbatch jobs/slurm/gpu.sh benchmark --operator fno      --run-name final_fno_seed$s
done
\end{cmd}

\begin{cmd}[Sunbird: when those have finished]
sbatch jobs/slurm/cpu.sh compare
\end{cmd}

\textbf{Stopping and continuing.} Jobs end at 24 hours. Submit the same \code{tune},
\code{train-scaling} or \code{generate} line again, or the same \code{train}/\code{train-final}
line with \code{--resume} at the end.

\textbf{Bringing results home.} In Anaconda Prompt on your computer, in the folder where the
results should go:

\begin{cmd}[Your computer]
scp -r USERNAME@LOGIN_NODE:/scratch/USERNAME/ML/outputs .
\end{cmd}

% =============================================================================
\part{Reference}
% =============================================================================

% =============================================================================
\section{All stages}
% =============================================================================

\begin{center}
\begin{tabularx}{\textwidth}{@{}llY@{}}
\toprule
Stage & Sunbird & What it does \\
\midrule
\code{generate}         & cpu & FE dataset (\code{--family scaling}: 3800 cases; default: 500) \\
\code{generate-rounded} & cpu & 50 rounded rectangles for an extra test (\code{--rounded-cases}) \\
\code{plot-sampling}    & cpu & figure of where the cavities lie \\
\code{preview}          & cpu & figure of one FE case: geometry, mesh, temperature, SDF \\
\code{mesh-check}       & cpu & FE accuracy on three meshes and heat balance (table) \\
\code{train}            & gpu & train one model \\
\code{evaluate}         & gpu & test errors of a run (\code{--split val} for validation) \\
\code{evaluate-rounded} & gpu & errors of a run on the rounded rectangles \\
\code{plot}             & gpu & loss curve and one test case (\code{--test-case N}) \\
\code{benchmark}        & gpu & training-epoch and prediction times \\
\code{tune}             & gpu & train the models of a tuning stage \\
\code{tune-select}      & cpu & pick the winner of a stage, draw its heatmap \\
\code{train-final}      & gpu & train the selected model with one seed \\
\code{summarize}        & cpu & mean and spread over the seeds of one operator \\
\code{compare}          & cpu & DeepONet against FNO: figures and tables \\
\code{train-scaling}    & gpu & data-size study: train on the first $N$ training cases \\
\code{scaling-report}   & cpu & data-size study: learning curves and rates $\alpha$ \\
\code{gpu-info}         & gpu & hardware record \\
\code{smoke}            & either & every stage on 20 tiny cases (a check) \\
\bottomrule
\end{tabularx}
\end{center}

% =============================================================================
\section{All options}\label{sec:options}
% =============================================================================

Every option has a default; type it only to change it. ``Full study'' marks the options the
study of Section~\ref{sec:order} uses.

\begin{center}
\small
\begin{tabularx}{\textwidth}{@{}llY@{}}
\toprule
Option & Default & Meaning \\
\midrule
\multicolumn{3}{@{}l}{\textit{Which data}} \\
\code{--family}          & \code{combined} & \code{scaling} (full study), \code{combined} (500-case study),
                                              \code{circle}, \code{rect} \\
\code{--cases}           & 500 / 3200      & cases to generate; for \code{scaling}: size of the training pool \\
\code{--train-size}      & 0 (all)         & use only the first $N$ training cases (full study: 800 for tuning) \\
\code{--sizes}           & 200,\ldots,3200 & sizes $N$ of \code{train-scaling} and \code{scaling-report} \\
\code{--val-cases}       & 200             & scaling dataset: validation cases \\
\code{--test-cases}      & 400             & scaling dataset: test cases \\
\code{--rounded-cases}   & 50              & rounded rectangles of \code{generate-rounded} \\
\code{--force}           & off             & make the dataset again although it exists \\
\code{--small-grid}      & off             & coarse grid and few cases (used by \code{smoke}) \\
\midrule
\multicolumn{3}{@{}l}{\textit{Which model and how it is trained}} \\
\code{--operator}        & \code{deeponet} & \code{deeponet} or \code{fno} (full study: both) \\
\code{--variant}         & \code{baseline} & named architecture of \code{train} (Section~\ref{sec:pcone}) \\
\code{--epochs}          & 5000            & epoch limit (full study: 2000 for tuning) \\
\code{--batch-size}      & 8               & cases per optimiser step \\
\code{--lr}              & 0.001           & initial learning rate (tuning overrides it in stage \code{lr}) \\
\code{--seed}            & 123             & seed of one run (initial weights, batch order) \\
\code{--seeds}           & 5 seeds         & 321,777,2026,31415,27182 for \code{train-scaling},
                                              \code{scaling-report}, \code{summarize}, \code{compare};
                                              123 for \code{tune} \\
\code{--run-name}        & automatic       & name of a run; all its files carry it \\
\code{--device}          & \code{cuda}     & \code{cuda} (GPU) or \code{cpu} \\
\code{--point-chunk}     & 2048            & DeepONet: grid points per pass (lower if GPU memory runs out;
                                              results unchanged) \\
\code{--cpu-threads}     & 4               & processor threads \\
\code{--checkpoint-every}& 50              & epochs between recovery points \\
\code{--early-stop}      & 800             & stop after this many epochs without a better validation error \\
\code{--scheduler-patience} & 200          & halve the learning rate after this many such epochs \\
\code{--resume}          & off             & continue a stopped \code{train}/\code{train-final} \\
\midrule
\multicolumn{3}{@{}l}{\textit{Evaluation, tuning, timing}} \\
\code{--split}           & \code{test}     & cases scored by \code{evaluate}: \code{test} or \code{val} \\
\code{--test-case}       & 0               & case drawn by \code{plot} \\
\code{--stage-name}      & --              & tuning stage (full study: every stage of Section~\ref{sec:pctune}) \\
\code{--cell}            & all             & one model of a tuning stage (on Sunbird: the array task) \\
\code{--benchmark-repeats} & 30            & timed repeats of a prediction \\
\code{--benchmark-epochs}  & 3             & timed training epochs \\
\code{--warmups}         & 5               & untimed repeats before timing \\
\bottomrule
\end{tabularx}
\end{center}

% =============================================================================
\section{Settings fixed in the code}\label{sec:settings}
% =============================================================================

These have no option on purpose: they define the problem, and changing one makes a new
dataset. They are all in \code{src/config.py}, each with a comment. To change one, edit the
number there and generate the data again (\code{--force}).

\begin{center}
\small
\begin{tabularx}{\textwidth}{@{}lY@{}}
\toprule
Group & Setting (value) \\
\midrule
Plate        & $L_x = 36$\,mm, $L_y = 21$\,mm \\
Physics      & $k = 28.3$\,W/(m\,K), $h = 4000$\,W/(m$^2$\,K), $T_\infty = 300$\,K,
               $q = 0.5$\,MW/m$^2$ on the top edge; other edges adiabatic \\
Cavities     & circle diameter 8--14\,mm; rectangle width 8--14\,mm, height 6--12\,mm;
               rounded corners 0.5--3\,mm (rounded set only) \\
FE           & H1 elements of order 2, element size $\le 0.8$\,mm \\
Grid         & $128 \times 96$ points; network input: SDF on this grid \\
500-case set & Latin hypercube, seed 42; split 80/10/10 per shape, seed 123 \\
Scaling set  & Halton pool of 3200; validation 200 and test 400 (Latin hypercube);
               sizes 200, 400, 800, 1600, 3200 \\
DeepONet     & branch 256--256, trunk 128--128, latent 128, tanh (baseline; tuning replaces depth and width) \\
FNO          & modes $12 \times 16$, width 32, 4 blocks, padding 8, GELU (baseline; tuning replaces them) \\
Training     & Adam; learning-rate factor 0.5; gradient clipping 1.0; 32-bit floats; deterministic \\
Tuning grids & in \code{src/training/tuning.py} (Section~\ref{sec:pctune}) \\
\bottomrule
\end{tabularx}
\end{center}

% =============================================================================
\section{When something goes wrong}
% =============================================================================

\begin{center}
\begin{tabularx}{\textwidth}{@{}>{\raggedright\arraybackslash}p{0.38\textwidth}Y@{}}
\toprule
The screen says & Remedy \\
\midrule
\code{Terms of Service have not been accepted} & create the environment with \code{--override-channels -c conda-forge} (Section~\ref{sec:env}) \\
\code{EnvironmentNameNotFound: opl}     & the environment was not created; repeat Section~\ref{sec:env} \\
\code{invalid choice: 'smoke--device'}  & a space is missing: \code{smoke --device cpu} \\
\code{CUDA is unavailable}              & CPU PyTorch is installed (Section~\ref{sec:env}), or on Sunbird the stage was sent to \code{cpu.sh} \\
\code{No module named ...}              & \code{conda activate opl} was forgotten, or a package is missing \\
\code{Run exists}                       & add \code{--resume}, or choose a new \code{--run-name} \\
\code{has not been selected yet}        & run \code{tune-select} of the previous stage first \\
\code{Cell not finished}                & some tuning models are still running or failed \\
\code{Checkpoint and dataset}           & the dataset changed after training; train again \\
\code{CUDA out of memory}               & add \code{--point-chunk 1024} or \code{--batch-size 4} \\
\code{Segmentation fault}               & \code{ngsolve} and \code{torch} used in one Python session; keep them apart \\
an error about \code{outputs/slurm}     & Sunbird: \code{mkdir -p outputs/slurm} in the project folder \\
\code{Connection timed out} (ssh, scp)  & your computer cannot reach Sunbird: connect to the university VPN first (or work on campus), then repeat \\
\bottomrule
\end{tabularx}
\end{center}

% =============================================================================
\section{One-page summary}
% =============================================================================

\begin{cmd}[Your computer (Anaconda Prompt): the full study]
conda activate opl
cd /d C:\ML

rem  steps 1-2: data
python run.py generate --family scaling
python run.py generate-rounded --family scaling --rounded-cases 50
python run.py plot-sampling --family scaling

rem  step 4: tuning at N = 800 (then lr; FNO: modes_width, depth, lr)
set T=--family scaling --train-size 800 --operator deeponet
python run.py tune        %T% --stage-name grid --epochs 2000
python run.py tune-select %T% --stage-name grid
python run.py tune        %T% --stage-name grid_gelu --epochs 2000
python run.py tune-select %T% --stage-name grid_gelu

rem  step 5: 25 runs per operator
python run.py train-scaling --operator deeponet
python run.py train-scaling --operator fno

rem  steps 6-7: results
python run.py scaling-report
for %n in (200 400 800 1600 3200) do python run.py compare --family scaling --train-size %n
\end{cmd}

\begin{cmd}[Sunbird (in the project folder)]
sbatch jobs/slurm/cpu.sh generate --family scaling
T="--family scaling --train-size 800 --operator fno"
A="--output=outputs/slurm/gpu_%A_%a.out"
sbatch --array=0-11 $A jobs/slurm/gpu.sh tune $T --stage-name modes_width --epochs 2000
sbatch jobs/slurm/cpu.sh tune-select $T --stage-name modes_width
sbatch --array=0-24 $A jobs/slurm/gpu.sh train-scaling --operator fno
sbatch jobs/slurm/cpu.sh scaling-report

squeue -u $USER
tail -f outputs/slurm/gpu_123456.out
scancel 123456
sacct -j 123456
\end{cmd}

\end{document}
